% Generated draft: requires mathematical and source review.
\documentclass[11pt]{article}
\usepackage[margin=25mm]{geometry}
\usepackage{fontspec}
\setmainfont{Latin Modern Roman}
\usepackage{amsmath,amssymb,mathtools,unicode-math}
\setmathfont{Latin Modern Math}
\usepackage{enumitem,needspace,xurl,hyperref}
\setlength{\parindent}{0pt}
\setlength{\parskip}{5pt}
\setlength{\emergencystretch}{3em}
\newcommand{\N}{\mathbb{N}}
\newcommand{\Z}{\mathbb{Z}}
\newcommand{\Q}{\mathbb{Q}}
\newcommand{\R}{\mathbb{R}}
\newcommand{\C}{\mathbb{C}}
\newcommand{\F}{\mathbb{F}}
\newcommand{\A}{\mathbb{A}}
\newcommand{\PP}{\mathbb{P}}
\newcommand{\E}{\mathbb{E}}
\newcommand{\eps}{\varepsilon}
\newcommand{\dd}{\,\mathrm{d}}
\begin{document}
\section*{UnsolvedMath ID: 800054}
\textbf{Source problem code:} catalogue entry
\subsection{Definitions and mathematical statement}
For an integer \(n\geq 0\), let \(\mathbb{P}^{n}(\R)=S^{n}/(x\sim -x)\) denote real projective space with its standard smooth structure. Define
\[
e(n)=\min\{m\in\Z_{\geq 0}:\text{there exists a smooth embedding }\mathbb{P}^{n}(\R)\hookrightarrow\R^{m}\},\qquad d(n)=e(n)-n.
\]
Here a smooth embedding is a smooth injective immersion that is a homeomorphism onto its image.

The supplied historical statement attributes to H. Hopf the bound \(d(n)>1\), including the assertion that this obstruction holds even for topological embeddings. It attributes to Chern the bound
\[
d(n)>2\quad\text{if}\quad n\notin\{2^{k}-1:k\in\Z,\ k\geq2\}\cup\{2^{k}-2:k\in\Z,\ k\geq2\}.
\]
These are reproduced as attributed source claims with an unresolved dimension qualification: the supplied record does not specify the range of \(n\), and the claims cannot hold as written for all positive integers, since \(\mathbb{P}^{1}(\R)\) is a circle and embeds smoothly in \(\R^{2}\). Thus the intended dimension range or low-dimensional exceptions must be checked before these bounds are adopted as hypotheses.

The general question is whether these lower bounds on the minimum embedding codimension can be improved, on their intended ranges of validity. The source does not prescribe a particular strengthened bound or a precise range for the improvement. Its explicit special question is:
\[
\text{Does there exist a smooth embedding }\mathbb{P}^{6}(\R)\hookrightarrow\R^{8}\text{?}
\]
Equivalently, is \(e(6)\leq8\), or \(d(6)\leq2\)? The smooth category here follows the source's definition of \(d(n)\); its separate topological qualification concerns the attributed Hopf obstruction.
\subsection{Short English statement}
Ask whether the historical lower bounds for the Euclidean embedding codimension of real projective space can be strengthened, and specifically whether real projective \(6\)-space embeds smoothly in \(\R^{8}\). The dimension range of the historical bounds needs clarification. The registry labels the problem open; this draft makes no independent assessment of its current status.
\subsection{Sources}
\begin{itemize}
\item UnsolvedMath local registry, permanent ID \(800054\), source problem code \texttt{catalogue entry}, title: Hirzebruch's 1954 Problem List --- Problem 3. The supplied registry records Kotschick as proposer and \(2013\) as proposed year, and labels the entry open; these are supplied metadata, not an independently verified historical attribution or current status assessment.
\item Kotschick, \emph{Updates on Hirzebruch's 1954 Problem List} (2013), Problem 3, supplied source URL: \url{https://arxiv.org/abs/1305.4623}. The background identifies the extraction method as \texttt{arxiv\_latex} and records access on 2026-10-16. This access date is supplied provenance; no independent source visit or literature verification is claimed.
\item The supplied statement attributes the first embedding obstruction to H. Hopf and the second to Chern. No separate bibliographic references for those results are supplied. The historical problem is identified as belonging to Hirzebruch's 1954 list, as reproduced in the cited 2013 update.
\end{itemize}
\end{document}
